\subsection{Справочник функций полёта}

Для управления полётом используется модуль \texttt{ap} (Autopilot). Ниже приведены основные функции, необходимые для выполнения миссий.

\begin{infobox}[Команды управления]
\begin{itemize}
   \item \texttt{ap.push(Ev.MCE\_PREFLIGHT)} — \textbf{Обязательная команда перед стартом.} Вводит дрон в режим предстартовой подготовки: проверяет датчики, калибрует гироскопы. Двигатели не запускаются.
   
   \item \texttt{ap.push(Ev.MCE\_TAKEOFF)} — Команда на взлёт. Дрон поднимается на высоту, заданную параметром \texttt{Flight\_com\_takeoffAlt} (обычно около 1 метра).
   
   \item \texttt{ap.push(Ev.MCE\_LANDING)} — Команда на посадку. Дрон плавно снижается, выключает двигатели после касания земли.
   
   \item \texttt{ap.goToLocalPoint(x, y, z)} — Полёт в точку в локальной системе координат.
   \begin{itemize}
       \item \texttt{x, y, z} — координаты в метрах.
       \item Ось X направлена вперёд, Y — влево, Z — вверх.
       \item Начало координат \((0,0,0)\) — фиксированная точка сцены (там, где дрон изначально размещён), а не место последнего взлёта: повторный вызов \texttt{ap.goToLocalPoint(0,0,z)} всегда возвращает дрон в одну и ту же точку, даже после нескольких перелётов.
   \end{itemize}
   
   \item \texttt{ap.updateYaw(rad)} — Поворот дрона вокруг своей оси (рысканье). Угол задаётся в радианах.
\end{itemize}
\end{infobox}

\begin{warnbox}[Аварийное выключение]
\texttt{ap.push(Ev.ENGINES\_DISARM)} — мгновенное выключение двигателей. Дрон упадёт камнем вниз.
\textbf{Использовать только в экстренных случаях} (например, угроза столкновения с человеком) или после гарантированной посадки.
\end{warnbox}

\subsubsection{Таймеры и события}
\begin{itemize}
   \item \texttt{Timer.callLater(delay, fn)} — «Сделать один раз через \texttt{delay} секунд».
   \item \texttt{Timer.new(period, fn)} — «Делать каждые \texttt{period} секунд».
   \item \texttt{function callback(event)} — Функция, которую система вызывает сама, когда что-то происходит (взлёт завершён, точка достигнута, села батарейка).
\end{itemize}
